\begin{song}{Jalava-Lied}{Schmetterlinge}

    \begin{verse}  
        Von [Am]Sonn' und Kessel schwarz gebrannt und auch vom scharfen [E]Wind, \\
        steht Jalava am Führerstand, wo Dampf und Flammen [Am]sind. \\
        Sein [Dm]neuer Heizer ist dabei, der [Am]ihm die Flamme nährt \\
        auf [E]Lokomotive zwei-neun-drei, die [Am]heut' nach Russland [E]fährt. \\
        Ein [Am]kleiner Mann von schmalem Bau, der werkt dort auf der [E]Brücke \\
        Ruß im Gesicht, das Haar war Grau, es war eine Pe[Am]rücke.
    \end{verse}  

    \begin{chorus}
        [Dm]Jalava, Jalava du Finne, was [Am]lachst du so gegen den Wind? \\
        Ich [E]lache, weil meine Sinne [Am]alle beisammen sind \\
        Und [Dm]weil wir weiterkamen und [Am]weil die Welt sich dreht \\
        Und [E]weil mein Heizer von Flammen und [Am]Dampfkesseln [E]was ver[Am]steht.
    \end{chorus}

    \begin{verse}  
        Sie [Am]dampfen ein in Belastrow, wo Schocks von Offi[E]zieren \\
        Die Züge auf dem Grenzbahnhof penibel kontrol[Am]lieren. \\
        Sie [Dm]prüfen jegliches Gesicht bei [Am]ihrer Inspizierung. \\
        Doch [E]sehen sie am Kessel nicht den [Am]Staatsfeind der Re[E]gierung. \\
        [Am]Jalava weiß worum es geht und [E]langsam dampft vorbei \\
        am letzten Posten der dort steht, Lokomo[Am]tive zwei-neun-drei.
    \end{verse}  

    \refchorus

    \textbf{Interlude:} [Dm Am E Am Dm Am E Am E Am]Jampa ja jalavala ...

    \begin{verse}  
        Jetzt [Am]saust die Grenzstation vorbei, die Birken stehen [E]nackt. \\
        Die Lokomotive zwei-neun-drei schnauft in erhöhtem [Am]Takt. \\
        Und [Dm]Jalava lacht in den Wind, in [Am]den Oktoberregen: \\
        "[E]Heizer, wenn wir drüben sind, dann wird sich was be[Am]we-[E]gen" \\
        Jetzt [Am]schneidet der Oktoberwind die letzten Äpfel [E]an, \\
        die an den kahlen Bäumen sind, an der finnischen Eisen[Am]bahn.
    \end{verse}  

    \begin{chorus}  
        [Dm]Jalava, Jalava du Finne, was [Am]lachst du so gegen den Wind? \\
        Ich [E]lache, weil meine Sinne [Am]alle beisammen sind \\
        Und [Dm]weil uns die Fahrt in den Bahnhof [Am]hinter die Grenze führt, \\
        und [E]Wladimir Iljitsch Uljanow, mein [Am]Heizer, die [E]Flammen [Am]schürt. \\
    \end{chorus}  

    \textbf{Outro:} [Dm Am E Am Dm Am E Am E Am]Jampa ja jalavala ...

\end{song}
